From - Mon May 22 18:45:08 2000
Received: from pcu.net (pcu.net [209.63.121.2])
	by math.usu.edu (8.9.3+Sun/8.9.1) with SMTP id OAA10247
	for <symanzik@sunfs.math.usu.edu>; Sat, 20 May 2000 14:40:55 -0600 (MDT)
Received: by pcu.net from localhost
    (router,SLMail V3.0); Sat, 20 May 2000 14:51:56 -0600
Received: by pcu.net from wzngemxm
    (216.190.144.86::mail daemon; unverified,SLMail V3.0); Sat, 20 May 2000 14:51:53 -0600
Message-ID: <000201bfc29c$43983b00$5690bed8@wzngemxm>
Reply-To: "Roxanne Drachenberg" <roxanned@pcu.net>
From: "Roxanne Drachenberg" <roxanned@pcu.net>
To: "Jurgen Symanzik" <symanzik@sunfs.math.usu.edu>
Subject: HW#6 - II
Date: Sat, 20 May 2000 14:39:37 -0600
MIME-Version: 1.0
Content-Type: multipart/alternative;
	boundary="----=_NextPart_000_0006_01BFC269.39245420"
X-Priority: 3
X-MSMail-Priority: Normal
X-Mailer: Microsoft Outlook Express 4.72.3155.0
X-MimeOLE: Produced By Microsoft MimeOLE V4.72.3155.0
Content-Length: 10491
Status: RO
X-Mozilla-Status: 8001
X-Mozilla-Status2: 00000000
X-UIDL: 38baa60a000002fa

This is a multi-part message in MIME format.

------=_NextPart_000_0006_01BFC269.39245420
Content-Type: text/plain;
	charset="iso-8859-1"
Content-Transfer-Encoding: quoted-printable

STAT 5810
Homework # 6
Part II

Evaluations
1)  Bruce and Glen, Describing data graphically
    They experienced some technical difficulties getting into the =
WebSite which got them off to a slow start.  I had trouble following =
along; took a while to get to the correct page.  Then Glen began his =
talk while Bruce was still navigating; it wasn't clear if we were to =
follow along the trail or listen to Glen--I lost track at that time.  =
Somehow the explanation on bins went over my head.  At one point, Glen =
was saying that you could manipulate the number of people in the study; =
but what he was really talking about was changing the number of bins =
which would change the length of the intervals and then affect the =
frequencies in each interval. This part of the presentation was not very =
effective.  Glen seemed hindered by having to use the web rather than =
the web being useful for him as a teaching tool.
    Bruce was able to use the web effectively however.  It is much =
quicker to be able to have a ready made box plot in pretty colors to =
point to and describe rather than having to draw one on the board.  =
However, making one on the board with the students working right along =
with you is very powerful for teaching HOW to make one.=20
    They had some nice support information in the form of handouts.  I =
think they just tried to get too much into a short lecture.

2)  Qian, JingJing, Summary Statistics
    They had great background information; gave great reasons why =
students should study statistics.  They refer to other sites for more =
information in the very beginning of their lecture.  They didn't wait =
for students to get on the same page.  I gave up and just listened =
rather than following along.  But this was fine for me. =20
    An overview of what they intended to accomplish was given in the =
beginning.  The examples to explain median and mean were great.  And =
they asked for questions to make sure we were understanding.  There was =
some confusion over the hands on experience example.  The web was used =
effectively.

3) Gary  Confidence Intervals
    I'm confused because Gary seemed to be evaluating a WebPage rather =
than teaching a course.  Maybe that's because he was absent earlier.  =
Anyway, since he isn't really a student of the class, I decline to =
evaluate his presentation.

4)  Weiping, Simple Linear Regression
    I thought this presentation made the most effective use of the web =
to present a lecture.  I found myself taking notes on the subject =
instead of evaluation notes on the presentation itself; that's how =
caught up in the class I became.  Her presentation was interesting, she =
gave a real world example for us to work through with her to enhance our =
understanding.  And at the end, she presented the formula but didn't bog =
us down by going through it.  I thought it was great.

    As far as which lecture was more effective than the pure "chalk and =
textbook" approach used in statistics classes today, I have to say that =
in my opinion, the effectiveness of a lecture still rests with the =
teacher, not just with the tools.  The most effective lecturer was =
Weiping; however, is this because she used the web or because she is in =
general a more effective teacher?  I think the web is a useful tool for =
the classroom, but as long as classes are still being taught in =
classrooms, I don't see it as the only tool.  A teacher who uses =
everything available to her and is a good teacher will be effective.

Personal Note:  I object to my comments on other student's presentations =
being made public whether or not my name accompanies them.  If I have =
any choice in the matter, I would prefer it if they were not dispersed =
among the class.  In fact, I don't want to read anyone else's =
impressions of my lecture either. Our comments are being graded and =
therefore should not be made public.  Thank you.
=20
Roxanne Drachenberg
(435) 752-8375
roxanned@pcu.net

------=_NextPart_000_0006_01BFC269.39245420
Content-Type: text/html;
	charset="iso-8859-1"
Content-Transfer-Encoding: quoted-printable

<!DOCTYPE HTML PUBLIC "-//W3C//DTD W3 HTML//EN">
<HTML>
<HEAD>

<META content=3Dtext/html;charset=3Diso-8859-1 =
http-equiv=3DContent-Type>
<META content=3D'"MSHTML 4.72.3110.7"' name=3DGENERATOR>
</HEAD>
<BODY bgColor=3D#ffffff>
<DIV><FONT color=3D#000000 size=3D2>STAT 5810</FONT></DIV>
<DIV><FONT color=3D#000000 size=3D2></FONT><FONT size=3D2>Homework # =
6</FONT></DIV>
<DIV><FONT size=3D2>Part II</FONT></DIV>
<DIV><FONT size=3D2></FONT>&nbsp;</DIV>
<DIV><FONT color=3D#000000 size=3D2>Evaluations</FONT></DIV>
<DIV><FONT color=3D#000000 size=3D2></FONT><FONT size=3D2>1)&nbsp; Bruce =
and Glen,=20
Describing data graphically</FONT></DIV>
<DIV><FONT size=3D2>&nbsp;&nbsp;&nbsp; They experienced some technical=20
difficulties getting into the WebSite which got them off to a slow =
start.&nbsp;=20
I had trouble following along; took a while to get to the correct =
page.&nbsp;=20
Then Glen began his talk while Bruce was still navigating; it wasn't =
clear if we=20
were to follow along the trail or listen to Glen--I lost track at that=20
time.&nbsp; Somehow the explanation on bins went over my head.&nbsp; At =
one=20
point, Glen was saying that you could manipulate the number of people in =
the=20
study; but what he was really talking about was changing the number of =
bins=20
which would change the length of the intervals and then affect the =
frequencies=20
in each interval. This part of the presentation was not very =
effective.&nbsp;=20
Glen seemed hindered by having to use the web rather than the web being =
useful=20
for him as a teaching tool.</FONT></DIV>
<DIV><FONT size=3D2></FONT><FONT color=3D#000000 =
size=3D2>&nbsp;&nbsp;&nbsp; Bruce was=20
able to use the web effectively however.&nbsp; It is much quicker to be =
able to=20
have a ready made box plot in pretty colors to point to and describe =
rather than=20
having to draw one on the board.&nbsp; However, making one on the board =
with the=20
students working right along with you is very powerful for teaching HOW =
to make=20
one.&nbsp;</FONT></DIV>
<DIV><FONT color=3D#000000 size=3D2>&nbsp;&nbsp;&nbsp; <FONT =
color=3D#000000>They had=20
some nice support information in the form of handouts.&nbsp; I think =
they just=20
tried to get too much into a short lecture.</FONT></FONT></DIV>
<DIV><FONT color=3D#000000 size=3D2><FONT =
color=3D#000000></FONT></FONT>&nbsp;</DIV>
<DIV><FONT size=3D2>2)&nbsp; Qian, JingJing, Summary =
Statistics</FONT></DIV>
<DIV><FONT size=3D2></FONT><FONT color=3D#000000 =
size=3D2>&nbsp;&nbsp;&nbsp; They had=20
great background information; gave great reasons why students should =
study=20
statistics.&nbsp; They refer to other sites for more information in the =
very=20
beginning of their lecture.&nbsp; They didn't wait for students to get =
on the=20
same page.&nbsp; I gave up and just listened rather than following =
along.&nbsp;=20
But this was fine for me.&nbsp; </FONT></DIV>
<DIV><FONT color=3D#000000 size=3D2>&nbsp;&nbsp;&nbsp; An overview of =
what they=20
intended to accomplish was given in the beginning.&nbsp; The examples to =
explain=20
median and mean were great.&nbsp; And they asked for questions to make =
sure we=20
were understanding.&nbsp; There was some confusion over the hands on =
experience=20
example.&nbsp; The web was used effectively.</FONT></DIV>
<DIV><FONT color=3D#000000 size=3D2></FONT>&nbsp;</DIV>
<DIV><FONT color=3D#000000 size=3D2>3) Gary&nbsp; Confidence =
Intervals</FONT></DIV>
<DIV><FONT color=3D#000000 size=3D2>&nbsp;&nbsp;&nbsp; I'm confused =
because Gary=20
seemed to be evaluating a WebPage rather than teaching a course.&nbsp; =
Maybe=20
that's because he was absent earlier.&nbsp; Anyway, since he isn't =
really a=20
student of the class, I decline to evaluate his =
presentation.</FONT></DIV>
<DIV><FONT color=3D#000000 size=3D2></FONT>&nbsp;</DIV>
<DIV><FONT color=3D#000000 size=3D2>4)&nbsp; Weiping, Simple Linear=20
Regression</FONT></DIV>
<DIV><FONT color=3D#000000 size=3D2>&nbsp;&nbsp;&nbsp; I thought this =
presentation=20
made the most effective use of the web to present a lecture.&nbsp; I =
found=20
myself taking notes on the subject instead of evaluation notes on the=20
presentation itself; that's how caught up in the class I became.&nbsp; =
Her=20
presentation was interesting, she gave a real world example for us to =
work=20
through with her to enhance our understanding.&nbsp; And at the end, she =

presented the formula but didn't bog us down by going through it.&nbsp; =
I=20
thought it was great.</FONT></DIV>
<DIV><FONT color=3D#000000 size=3D2></FONT>&nbsp;</DIV>
<DIV><FONT color=3D#000000 size=3D2>&nbsp;&nbsp;&nbsp; As far as which =
lecture was=20
more effective than the pure &quot;chalk and textbook&quot; approach =
used in=20
statistics classes today, I have to say that in my opinion, the =
effectiveness of=20
a lecture still rests with the teacher, not just with the tools.&nbsp; =
The most=20
effective lecturer was Weiping; however, is this because she used the =
web or=20
because she is in general a more effective teacher?&nbsp; I think the =
web is a=20
useful tool for the classroom, but as long as classes are still being =
taught in=20
classrooms, I don't see it as the only tool.&nbsp; A teacher who uses =
everything=20
available to her and is a good teacher will be effective.</FONT></DIV>
<DIV><FONT color=3D#000000 size=3D2></FONT>&nbsp;</DIV>
<DIV><FONT color=3D#000000 size=3D2><STRONG><U>Personal =
Note</U></STRONG>:&nbsp;=20
<EM>I object to my comments on other student's presentations being made =
public=20
whether or not my name accompanies them.&nbsp; If I have any choice in =
the=20
matter, I would prefer it if they were not dispersed among the =
class.&nbsp; In=20
fact, I don't want to read anyone else's impressions of my lecture =
either. Our=20
comments are being graded and therefore should not be made public.&nbsp; =
Thank=20
you.</EM></FONT></DIV>
<DIV><FONT color=3D#000000 size=3D2><EM>&nbsp;</EM></FONT></DIV>
<DIV><FONT color=3D#000000 size=3D2>Roxanne Drachenberg<BR>(435) =
752-8375<BR><A=20
href=3D"mailto:roxanned@pcu.net">roxanned@pcu.net</A></FONT></DIV></BODY>=
</HTML>

------=_NextPart_000_0006_01BFC269.39245420--

